\documentclass[10pt,a4paper]{amsart}
\usepackage[margin=19mm]{geometry}
\usepackage{amsmath,amssymb,mathtools}
\usepackage[scr=rsfs,cal=euler]{mathalfa}
\usepackage{xfrac}
\usepackage[unicode,hidelinks,bookmarks=false]{hyperref}
\newcommand{\ie}{i.e.\ }
\newcommand{\cf}{cf.\ }
\newcommand{\resp}{respectively\ }
\newcommand{\Subsection}{\subsection}
\providecommand{\nobreakdash}{\nobreak\hbox{-}}
\setlength{\emergencystretch}{2em}
\multlinegap0pt
\usepackage{bm}
\usepackage{bbm}
\usepackage{dsfont}
\usepackage{xspace}
\usepackage{cancel}
\usepackage{mathtools}
\usepackage[shortlabels]{enumitem}
\usepackage{amsthm}
\makeatletter
\@ifundefined{theorem}{\newtheorem{theorem}{Theorem}}{}
\@ifundefined{lemma}{\newtheorem{lemma}[theorem]{Lemma}}{}
\makeatother
\title[Stochastic realisers of non-degenerate full-degree Type II r]{Stochastic realisers of non-degenerate full-degree Type II reduced Ito polynomials}
\author{Brecht Verbeken, Vincent Ginis}
\date{Source: arXiv:2609.26055v1 (CC BY 4.0)}
\begin{document}
\maketitle

\noindent\emph{Source and licence.} Stochastic realisers of non-degenerate full-degree Type II reduced Ito polynomials. Version arXiv:2609.26055v1 (CC BY 4.0), distributed under the Creative Commons Attribution 4.0 International licence, as recorded in the arXiv licence metadata for this exact version and shown on the abstract page https://arxiv.org/abs/2609.26055v1. Full rights notice: licenses/arxiv-2609.26055-CC-BY-4.0.txt.\par\smallskip

\noindent\textit{Editorial scope.} Counted unit: the Theorem whose source label is \texttt{thm:support} in the article (source0.tex lines 702--723, source label \texttt{thm:support}), delivered together with the article's own complete original proof of it (source0.tex lines 725--1004, one \texttt{proof} environment of 280 lines). No mathematics is written, repaired or re-derived: statement and proof are the article's own text line for line. Required context retained uncounted: eq:horiz (lines 479--482); eq:transferlast (lines 490--493); lem:tuple-simple (lines 566--583); eq:lengthone (lines 577--582); eq:constantcarry (lines 592--597); lem:kone (lines 600--603); lem:separation (lines 633--657); eq:monge (lines 638--641). Every definition, display and label of those lines is retained unchanged. Omitted from this package and not redistributed: 1--478, 483--489, 494--565, 583--591, 598--599, 604--632, 642--701, 724--724, 1005--1379. Nothing inside a retained excerpt is omitted or altered: every retained body is delivered exactly as the article's lines are written, so no omission receipt of any kind accompanies this package. Citations: the retained excerpts carry no citation command at all, so no bibliography entry of the article is transcribed and no reference is redistributed. Reference adaptation: none was needed and none was made; every reference inside the delivered unit resolves inside it, so no unresolved reference or citation is produced. Printed slips kept literal: no printed word, symbol, hypothesis or number of the counted statement or of its proof was altered. Layout and class: the article's own class is \texttt{article} with packages including inputenc, fontenc, lmodern, microtype, amsmath, amssymb, amsthm, mathtools, graphicx, geometry, enumitem, hyperref; the wrapper is the lane's pinned \texttt{amsart} editorial wrapper at 10pt on A4, which restates the theorem-like environments the retained text needs and adds the article's own macro definitions that the retained text needs (none), with extra wrapper packages (bm, bbm, dsfont, xspace, cancel, mathtools, enumitem). The delivered numbering is the unit's own. Assets: no retained span contains a figure environment, a graphics-inclusion command, a PDF-page inclusion, a table or a TikZ picture; no image file is copied into this package, the package declares no asset dependency and no image file is redistributed with it. Licence: version arXiv:2609.26055v1 of this article is distributed under the Creative Commons Attribution 4.0 International licence, as recorded in the arXiv licence metadata for this exact version and shown on the abstract page https://arxiv.org/abs/2609.26055v1. A full rights notice is delivered at licenses/arxiv-2609.26055-CC-BY-4.0.txt. Characters: every retained excerpt is pure ASCII, so the delivered engine, XeLaTeX with its default Latin Modern text font, reports no missing character.
\begin{equation}
        h_{t,i}:(t,i)\longrightarrow(t,i+1),
        \label{eq:horiz}
\end{equation}

\begin{equation}
        g_{d-1,i}:(d-1,i)\longrightarrow(0,i+c).
        \label{eq:transferlast}
\end{equation}

\begin{lemma}[Transfer tuples are simple cycles]
\label{lem:tuple-simple}
Let $S_0,\ldots,S_{d-1}$ be nonempty transfer sets in the two-choice graph \eqref{eq:horiz}--\eqref{eq:transferlast}.  For every tuple
\[
        (i_0,\ldots,i_{d-1})\in S_0\times\cdots\times S_{d-1},
\]
there is a simple directed cycle using exactly the transfer edges
\[
        g_{0,i_0},g_{1,i_1},\ldots,g_{d-1,i_{d-1}}.
\]
Its length is
\begin{equation}
        L(i_0,\ldots,i_{d-1})
        =d+\sum_{t=1}^{d-1}\langle i_t-i_{t-1}\rangle_q
          +\langle i_0-i_{d-1}-c\rangle_q .
        \label{eq:lengthone}
\end{equation}
\end{lemma}

\begin{equation}
        \sum_{t=1}^{d-1}\langle i_t-i_{t-1}\rangle_q
          +\langle i_0-i_{d-1}-c\rangle_q
        =dq-c .
        \label{eq:constantcarry}
\end{equation}

\begin{lemma}[Non-horizontal cycles use one transfer edge per block]
\label{lem:kone}
If a Type II admissible support has a directed cycle using at least one transfer edge, and if it uses $k$ transfer edges from each block, then $k=1$.
\end{lemma}

\begin{lemma}[Circular Monge separation]
\label{lem:separation}
Let $A,B\subseteq\mathbb Z_q$ be nonempty.  The following are equivalent.
\begin{enumerate}[label=\textup{(\roman*)}]
\item For all $a,a'\in A$ and $b,b'\in B$,
\begin{equation}
        D(a,b)+D(a',b')=D(a,b')+D(a',b).
        \label{eq:monge}
\end{equation}
\item The family of directed arcs $[a,b]_q$, $a\in A$, $b\in B$, has a common blocker.
\item There are integer representatives $\widehat a\equiv a\pmod q$, $\widehat b\equiv b\pmod q$, chosen for all $a\in A$ and $b\in B$, such that
\begin{equation}
        \widehat b<\widehat a\le \widehat b+q
        \qquad(a\in A,\ b\in B),
        \label{eq:separated-lifts-order}
\end{equation}
 and hence
\begin{equation}
        D(a,b)=q+\widehat b-\widehat a
        \qquad(a\in A,\ b\in B).
        \label{eq:separated-lifts-distance}
\end{equation}
\item There are integer intervals $I_B=[u,v]$ and $I_A=[v+1,w]$, with $w-u\le q$, whose residue classes contain $B$ and $A$, respectively.  Equivalently, $B$ and $A$ are contained in consecutive cyclic intervals whose total cardinality is at most $q+1$, with the interval containing $B$ immediately preceding the interval containing $A$.
\end{enumerate}
\end{lemma}

\begin{theorem}[Type II support theorem]
\label{thm:support}
Let $q\ge2$, $d\ge2$, $1\le z\le q-1$, and $\gcd(q,z)=1$. Put $c=d+z$ and $s=qd-z$. Let $S_0,\ldots,S_{d-1}$ be nonempty subsets of $\mathbb Z_q$ in the Type II two-choice digraph \eqref{eq:horiz}--\eqref{eq:transferlast}. Then the support is Type II admissible if and only if there exist
\[
        a_0,\ldots,a_{d-1}\in\mathbb Z_{\ge0},
        \qquad a_0+\cdots+a_{d-1}=z,
\]
and $\rho\in\mathbb Z_q$ such that
\begin{equation}
        S_0\subseteq [\rho,\rho+a_0]_q,
        \label{eq:S0a}
\end{equation}
\begin{equation}
        S_t\subseteq
        \left[
        \rho-\sum_{j=1}^t(a_j+1),\,
        \rho-\sum_{j=1}^t(a_j+1)+a_t
        \right]_q,
        \qquad t=1,\ldots,d-1 .
        \label{eq:Sta}
\end{equation}
\end{theorem}

\begin{proof}
\noindent\emph{Sufficiency.}
Assume that the displayed composition and interval conditions hold, and put
\[
        m_t=a_t+1,
        \qquad t\in\mathbb Z_d.
\]
Then
\[
        \sum_{t=0}^{d-1}m_t=d+z=c,
        \qquad
        m_t\le z+1\le q,
        \qquad
        m_{t-1}+m_t=2+a_{t-1}+a_t\le 2+z\le q+1,
\]
so the bounds used below are automatic consequences of the composition. Choose integer lifts
\[
        R_0=\rho,
        \qquad
        R_t=\rho-\sum_{j=1}^t m_j\quad (t=1,\ldots,d-1),
\]
and
\[
        R_d=R_0-(d+z)=R_0-c .
\]
Then
\begin{equation}
        S_t\subseteq [R_t,R_{t-1}-1]_q,
        \qquad t=1,\ldots,d-1,
        \label{eq:suff-incl1}
\end{equation}
and, after translating the interval for $S_0$ by $-c$,
\begin{equation}
        S_0-c\subseteq [R_d,R_{d-1}-1]_q .
        \label{eq:suff-incl0}
\end{equation}
For each adjacent pair of transfer sets, the target set and the source set are contained in consecutive cyclic intervals of lengths $m_t$ and $m_{t-1}$, with indices read modulo $d$. By the preceding local bound and Lemma \ref{lem:separation}, all horizontal arcs in that block have a common blocker. Explicitly, this applies to the arcs from $S_{t-1}$ to $S_t$ for $t=1,\ldots,d-1$, and to the arcs from $S_{d-1}$ to $S_0-c$, equivalently from $S_{d-1}+c$ to $S_0$, in block $0$.

A simple directed cycle cannot use two transfer edges with target in the
same block.  Indeed, between the two transfer targets and the two subsequent
transfer tails, the cycle would contain two horizontal subpaths in that
block with disjoint vertex sets; Lemma \ref{lem:separation} gives a common
blocker for all such subpaths, a contradiction.  Since the number of
transfer edges used in each block is the same, every transfer cycle uses
exactly one transfer edge from each block.

It remains to compute its length.  Let $i_t\in S_t$.  Choose lifts so that
\[
        \widetilde i_0\in [R_0,R_0+m_0-1],
        \qquad
        \widetilde i_t\in [R_t,R_{t-1}-1]\quad(t=1,\ldots,d-1),
\]
so that $\widetilde i_0-c\in [R_d,R_{d-1}-1]$. The interval inclusions, together with the preceding local bound, give
\[
        1\le \widetilde i_{t-1}-\widetilde i_t
        \le m_{t-1}+m_t-1\le q,
        \qquad t=1,\ldots,d-1,
\]
and
\[
        1\le \widetilde i_{d-1}-(\widetilde i_0-c)
        \le m_{d-1}+m_0-1\le q.
\]
Hence
\[
        \langle i_t-i_{t-1}\rangle_q=q+\widetilde i_t-\widetilde i_{t-1},
        \qquad t=1,\ldots,d-1,
\]
and
\[
        \langle i_0-i_{d-1}-c\rangle_q=q+\widetilde i_0-c-\widetilde i_{d-1}.
\]
Substitution in \eqref{eq:lengthone} gives
\[
\begin{aligned}
        L
        &=d+\sum_{t=1}^{d-1}(q+\widetilde i_t-\widetilde i_{t-1})
          +(q+\widetilde i_0-c-\widetilde i_{d-1})  \\
        &=d+dq-c
        =qd-z .
\end{aligned}
\]
Thus every transfer cycle has length $s=qd-z$, and the only cycles without transfer edges are the $d$ horizontal $q$-cycles.  The support is admissible.

\medskip
\noindent\emph{Necessity.}
Assume the support is Type II admissible.  By Lemma \ref{lem:tuple-simple}, every tuple $(i_0,\ldots,i_{d-1})\in S_0\times\cdots\times S_{d-1}$ gives a simple directed transfer cycle using exactly one transfer edge from each block.  Since the support is admissible, every such cycle has length $s=qd-z$.  Consequently \eqref{eq:constantcarry} holds for every tuple.

It is convenient to put
\[
        T_t=S_t\quad(0\le t\le d-1),
        \qquad
        T_d=S_0-c,
\]
and, for $t=1,\ldots,d$, to write
\[
        E_t(x,y)=D(x,y)=\langle y-x\rangle_q,
        \qquad x\in T_{t-1},\ y\in T_t .
\]
Here an element $y\in T_d$ is always of the form $y=i_0-c$ with $i_0\in S_0$.  The constant-carry identity becomes
\begin{equation}
        \sum_{t=1}^d E_t(x_{t-1},x_t)=dq-c,
        \qquad x_t\in T_t,
        \qquad x_d=x_0-c .
        \label{eq:Tconstant}
\end{equation}

We first prove that every adjacent pair $(T_{t-1},T_t)$ satisfies the Monge identity \eqref{eq:monge}.  Suppose $d\ge3$.  Fix $t\in\{1,\ldots,d\}$, and choose
\[
        a,a'\in T_{t-1},
        \qquad
        b,b'\in T_t .
\]
We show that
\[
        E_t(a,b)+E_t(a',b')=E_t(a,b')+E_t(a',b).
\]
There are three cases, because the cyclic constraint $x_d=x_0-c$ makes the endpoints look slightly different.

If $2\le t\le d-1$, fix all variables except $x_{t-1}$ and $x_t$; in particular fix $u\in T_{t-2}$ and $v\in T_{t+1}$.  In \eqref{eq:Tconstant}, the part depending on $x_{t-1}=a$ and $x_t=b$ is
\[
        E_{t-1}(u,a)+E_t(a,b)+E_{t+1}(b,v)+C,
\]
where $C$ is independent of $a,b$.  Since the whole sum in \eqref{eq:Tconstant} is the same for the four choices $(a,b)$, $(a,b')$, $(a',b)$, $(a',b')$, adding the identities for $(a,b)$ and $(a',b')$ and subtracting the identities for $(a,b')$ and $(a',b)$ cancels the two neighbouring terms and $C$, leaving exactly the Monge identity for $E_t$.

For $t=1$, the variable $a=x_0\in T_0$ also determines the endpoint $x_d=a-c\in T_d$.  Fix choices $x_2^*,\ldots,x_{d-1}^*$, with the evident interpretation that for $d=3$ there is only the single fixed choice $x_2^*$.  The part of \eqref{eq:Tconstant} depending on $a\in T_0$ and $b\in T_1$ is
\[
        E_d(x_{d-1}^*,a-c)+E_1(a,b)+E_2(b,x_2^*)+C .
\]
The same four-corner cancellation cancels $E_d(x_{d-1}^*,a-c)$, $E_2(b,x_2^*)$, and $C$, and gives the Monge identity for $E_1$.

For $t=d$, write $b\in T_d$ and remember that the corresponding initial variable is $x_0=b+c\in T_0$.  Fix choices $x_1^*,\ldots,x_{d-2}^*$, again with only one fixed choice when $d=3$.  The part depending on $a\in T_{d-1}$ and $b\in T_d$ is
\[
        E_{d-1}(x_{d-2}^*,a)+E_d(a,b)+E_1(b+c,x_1^*)+C .
\]
Again the four-corner cancellation gives the Monge identity for $E_d$.  Thus every adjacent pair satisfies \eqref{eq:monge} when $d\ge3$.

Now let $d=2$, and write $A=S_0$ and $B=S_1$.  For every $a\in A$ and $b\in B$, the one-transfer cycle has length $s=2q-z$, hence
\[
        D(a,b)+D(b+c,a)=2q-c .
\]
Equivalently,
\begin{equation}
        D(a,b)>q-c
        \qquad(a\in A,\ b\in B).
        \label{eq:d2-threshold}
\end{equation}
Suppose that the arcs $[a,b]_q$, $a\in A$, $b\in B$, have no common blocker.  By the contrapositive part of Lemma \ref{lem:separation}, there are two disjoint such arcs.  After relabelling the two arcs if necessary, rotate the circle and choose integer representatives so that
\[
        a=0\le b<a'\le b'<q .
\]
The inequalities \eqref{eq:d2-threshold} applied to $(a,b)$ and $(a',b')$ give
\[
        b+c>q,
        \qquad
        b'+c-q>a' .
\]
Also $c\le q+1$ and $b<a'$ give
\[
        b+c-q\le b+1\le a' .
\]
Thus in block $1$ the arcs $[a,b]$ and $[a',b']$ are disjoint, while in block $0$ the crossed arcs
\[
        [b+c-q,a']
        \qquad\text{and}\qquad
        [b'+c-q,q]
\]
are disjoint lifts of the arcs from $b+c$ to $a'$ and from $b'+c$ to $a$.  The corresponding cycle is, in lifted notation,
\[
\begin{aligned}
        (0,a)&\to(1,a)\leadsto(1,b)\to(0,b+c)
        \leadsto(0,a')  \\
        &\to(1,a')\leadsto(1,b')\to(0,b'+c)
        \leadsto(0,a),
\end{aligned}
\]
where $\leadsto$ denotes the appropriate horizontal path inside the indicated
block.  The transfer tails are distinct in their respective blocks, and the
transfer heads are precisely the initial vertices of the displayed horizontal
paths, so the transfer edges introduce no additional vertex identifications.
The two block-$1$ horizontal paths are disjoint by $a=0\le b<a'\le b'$, and the
two block-$0$ horizontal paths are disjoint by the displayed crossed intervals.
Hence this is a simple directed cycle using two transfer edges from each block,
contradicting Lemma \ref{lem:kone}.  Thus the arcs from $A$ to $B$ have a common
blocker, and Lemma \ref{lem:separation} gives the Monge identity for $(A,B)$.
The final pair satisfies the Monge identity as well, because
\[
        D(b,a-c)=2q-c-D(a,b)
        \qquad(a\in A,\ b\in B)
\]
by the one-transfer length condition; replacing $D(a,b)$ by this affine negative preserves the $2\times2$ identity.

By Lemma \ref{lem:separation}, choose separated integer lifts $X_0:T_0\to\mathbb Z$ and $X_1:T_1\to\mathbb Z$ for the first adjacent pair, so that
\begin{equation}
        E_1(x_0,x_1)=q+X_1(x_1)-X_0(x_0)
        \qquad(x_0\in T_0,
        \ x_1\in T_1).
        \label{eq:lift1}
\end{equation}
We now propagate these lifts around the cycle.  Suppose $2\le t\le d$ and $X_{t-1}$ has already been defined so that
\[
        E_{t-1}(x_{t-2},x_{t-1})
        =q+X_{t-1}(x_{t-1})-X_{t-2}(x_{t-2}).
\]
Fix $x_{t-2}\in T_{t-2}$ and $x_t\in T_t$ and vary $x_{t-1}\in T_{t-1}$ in \eqref{eq:Tconstant}, keeping all other coordinates fixed.  The terms depending on $x_{t-1}$ are
\[
        q+X_{t-1}(x_{t-1})-X_{t-2}(x_{t-2})
        +E_t(x_{t-1},x_t),
\]
so $X_{t-1}(x_{t-1})+E_t(x_{t-1},x_t)$ is independent of $x_{t-1}$.  Choosing a base point $x_{t-1}^*\in T_{t-1}$ and setting
\[
        X_t(x_t)=X_{t-1}(x_{t-1}^*)+E_t(x_{t-1}^*,x_t)-q
\]
therefore gives
\begin{equation}
        E_t(x_{t-1},x_t)=q+X_t(x_t)-X_{t-1}(x_{t-1})
        \qquad(x_{t-1}\in T_{t-1},
        \ x_t\in T_t).
        \label{eq:liftt}
\end{equation}
This defines compatible lifts $X_t$ for all $t=0,1,\ldots,d$.  Moreover, each $X_t$ is an integer representative of the corresponding residue class:
\begin{equation}
        X_t(x)\equiv x\pmod q,
        \qquad x\in T_t .
        \label{eq:lift-residue}
\end{equation}
This is true for $X_0$ and $X_1$ by the initial choice of separated lifts.  If it is true for $X_{t-1}$, then the defining formula for $X_t$ gives
\[
        X_t(x_t)\equiv x_{t-1}^*+\langle x_t-x_{t-1}^*\rangle_q\equiv x_t\pmod q,
\]
so \eqref{eq:lift-residue} follows by induction.

Summing \eqref{eq:liftt} over $t=1,\ldots,d$ and comparing with \eqref{eq:Tconstant} gives, for every $x_0\in T_0$,
\[
        dq+X_d(x_0-c)-X_0(x_0)=dq-c .
\]
Hence
\begin{equation}
        X_d(x_0-c)=X_0(x_0)-c
        \qquad(x_0\in S_0).
        \label{eq:closure-lifts}
\end{equation}
Now define
\[
        R_t=\min X_t(T_t),
        \qquad t=0,1,\ldots,d .
\]
By \eqref{eq:closure-lifts}, $R_d=R_0-c$.  Moreover, \eqref{eq:liftt} and $0\le E_t\le q-1$ imply
\[
        1\le X_{t-1}(x_{t-1})-X_t(x_t)\le q
        \qquad(x_{t-1}\in T_{t-1},
        \ x_t\in T_t).
\]
Since the $X_t$ are integer representatives by \eqref{eq:lift-residue}, these inequalities imply
\begin{equation}
        T_t\subseteq [R_t,R_{t-1}-1]_q,
        \qquad t=1,\ldots,d .
        \label{eq:nec-lift-inclusion}
\end{equation}
For $t=d$ this says $S_0-c\subseteq [R_d,R_{d-1}-1]_q$; adding $c$ and using $R_d+c=R_0$ gives the required interval for $S_0$.

Define
\[
        m_t=R_{t-1}-R_t,
        \qquad t=1,\ldots,d-1,
\]
and
\[
        m_0=R_{d-1}-R_d .
\]
The preceding inequalities give $m_t\in\{1,\ldots,q\}$.  Also
\[
        \sum_{t=0}^{d-1}m_t=R_0-R_d=c=d+z .
\]
Consequently, writing $a_t=m_t-1$, we have $a_t\ge0$ and $\sum_t a_t=z$. Moreover,
\[
        m_{t-1}+m_t=2+a_{t-1}+a_t\le 2+z\le q+1 .
\]
The inclusions \eqref{eq:nec-lift-inclusion}, together with the translated inclusion for $S_0$, are exactly \eqref{eq:S0a}--\eqref{eq:Sta} with $\rho=R_0$. This completes the proof.
\end{proof}

\end{document}
